\documentclass[11pt,a4paper]{article}
\usepackage[T1]{fontenc}
\usepackage[utf8]{inputenc}
\usepackage{amsmath,amssymb,xcolor}
\usepackage[margin=25mm]{geometry}
\usepackage[colorlinks=true,urlcolor=blue]{hyperref}
\definecolor{nexorablue}{HTML}{2459B5}
\setlength{\parindent}{0pt}
\setlength{\parskip}{7pt}
\setlength{\emergencystretch}{2em}
\renewcommand{\familydefault}{\sfdefault}
\begin{document}

{\small\bfseries\color{nexorablue} NEXORA CPGE}\par
{\small MP / NEX-RED-04 / Corrigé}\par
{\LARGE\bfseries Calculer des puissances sur deux directions stables}\par
\vspace{0.25em}\hrule\vspace{0.6em}

\section*{Directions propres}
Les vecteurs $(1,1)$ et $(1,-1)$ sont propres, pour $3$ et $1$. Ils sont indépendants. Un vecteur $(x,y)$ s’écrit $(x+y)(1,1)/2+(x-y)(1,-1)/2$.

\section*{Puissance}
La première composante propre est multipliée par $3^n$, la seconde par $1$. En appliquant la formule aux deux vecteurs de la base canonique, on obtient
\[ A^n=\frac12\begin{pmatrix}3^n+1&3^n-1\\3^n-1&3^n+1\end{pmatrix}. \]


\end{document}
